\SourceLeaf{285}%
\begin{center}\begin{tabular}{r|ccc}
Starting point&&&\\
&$-a$&$-b$&$+c$\\
8&&&\\
&$-a$&&\\
9&&&\\
&$+a$&&\\
12&&&\\
&$-a$&&\\
14&&&\\
&$+a$&&\\
16&&&\\
&$+a$&$-b$&\\
17&&&\\
&$-a$&$-b$&\\
18&&&\\
&$+a$&$-b$&\\
19&&&\\
&$+a$&$-b$&$+c$\\
20&&&\\
&$+a$&$+b$&$+c$\\
21&&&\\
&$-a$&$+b$&$+c$\\
22&&&\\
&$-a$&$-b$&$+c$\\
Starting point&&&\\
\end{tabular}\end{center}
The preceding analysis shows that the last sign corresponding
to each figure in this table must necessarily be the same
as the first. In other words, the signs $+-$ occurring at the individual
intersections cancel one another if we set aside those occurring at the
individual points where the figures form.
The sum of \textls{all} the signs is therefore the same as the sum of the latter alone, and
is also identical to the sum of the signs in the character of the
starting point relative to the individual figures $a$, $b$, $c$, i.\,e. all of them
if one does not count the last to be formed, $d$.
\par\Needspace{12\baselineskip}
Let the character on one side of the figure $d$ relative to the individual
figures be ex[pressed] by $\alpha a+\beta b+\gamma c+\delta d$
and [let us] set $\alpha+\beta+\gamma+\delta=\sigma$. Let the absolute
character of that location, on the other hand, be $=\chi$. The total
lateral turning is then
\[
= -\sigma\,.\,360^0 - \frac{F-\chi S}{S}\,.\,720^0,
\]
